\documentclass[UTF8,12pt]{ctexart}
\usepackage[a4paper,margin=24mm]{geometry}
\usepackage{amsmath,amssymb,bm,graphicx,adjustbox}
\newcommand{\fitmath}[1]{\begin{adjustbox}{max width=\linewidth}$\displaystyle #1$\end{adjustbox}}
\setlength{\parindent}{0pt}
\setlength{\parskip}{.7em}
\sloppy
\allowdisplaybreaks
\begin{document}
\section*{1999 年考研数学三 · 真题}
\subsection*{1999 年全国硕士研究生招生考试试题}

\subsection*{一、填空题（本题共 5 小题，每小题 3 分，满分 15 分）}

（1）设 \(\displaystyle f(x)\) 有一个原函数 \(\displaystyle \dfrac{\displaystyle \sin x}{\displaystyle x}\)，则 \(\displaystyle \int_{\pi/2}^{\pi}xf'(x)\,\mathrm dx=\)\_\_\_\_\_\_。


（2）\(\displaystyle \sum_{n=1}^{\infty}n\left(\dfrac{\displaystyle 1}{\displaystyle 2}\right)^{n-1}=\)\_\_\_\_\_\_。


（3）设 \(\displaystyle A=\begin{pmatrix}1&0&1\\0&2&0\\1&0&1\end{pmatrix}\)，而 \(\displaystyle n\ge2\) 为正整数，则 \(\displaystyle A^n-2A^{n-1}=\)\_\_\_\_\_\_。


（4）在天平上重复称量一重为 \(\displaystyle a\) 的物品，假设各次称量结果相互独立且同服从正态分布 \(\displaystyle N(a,\allowbreak 0.2^2)\)。若以 \(\displaystyle \bar X_n\) 表示 \(\displaystyle n\) 次称量结果的算术平均值，则为使 \(\displaystyle P\{|\bar X_n-a|<0.1\}\ge0.95\)，\(\displaystyle n\) 的最小值应不小于自然数\_\_\_\_\_\_。


（5）设随机变量 \(\displaystyle X_{ij}\;(i,\allowbreak j=1,\allowbreak 2,\allowbreak \ldots,\allowbreak n;\ n\ge2)\) 独立同分布，\(\displaystyle E(X_{ij})=2\)，则行列式 \(\displaystyle Y=\begin{vmatrix}X_{11}&X_{12}&\cdots&X_{1n}\\X_{21}&X_{22}&\cdots&X_{2n}\\\vdots&\vdots&&\vdots\\X_{n1}&X_{n2}&\cdots&X_{nn}\end{vmatrix}\) 的数学期望 \(\displaystyle E(Y)=\)\_\_\_\_\_\_。


\subsection*{二、选择题（本题共 5 小题，每小题 3 分，满分 15 分）}

（1）设 \(\displaystyle f(x)\) 是连续函数，\(\displaystyle F(x)\) 是 \(\displaystyle f(x)\) 的原函数，则（　）。（A）当 \(\displaystyle f(x)\) 是奇函数时，\(\displaystyle F(x)\) 必为偶函数。（B）当 \(\displaystyle f(x)\) 是偶函数时，\(\displaystyle F(x)\) 必为奇函数。（C）当 \(\displaystyle f(x)\) 是周期函数时，\(\displaystyle F(x)\) 必为周期函数。（D）当 \(\displaystyle f(x)\) 是单调增函数时，\(\displaystyle F(x)\) 必为单调增函数。


（2）设 \(\displaystyle f(x,\allowbreak y)\) 连续，且 \(\displaystyle f(x,\allowbreak y)=xy+\displaystyle \iint_D f(u,\allowbreak v)\,\mathrm du\mathrm dv\)，其中 \(\displaystyle D\) 是由 \(\displaystyle y=0,\allowbreak y=x^2,\allowbreak x=1\) 所围区域，则 \(\displaystyle f(x,\allowbreak y)\) 等于（　）。（A）\(\displaystyle xy\)。（B）\(\displaystyle 2xy\)。（C）\(\displaystyle xy+\dfrac{\displaystyle 1}{\displaystyle 8}\)。（D）\(\displaystyle xy+1\)。


（3）设向量 \(\displaystyle \beta\) 可由向量组 \(\displaystyle \alpha_1,\allowbreak \alpha_2,\allowbreak \ldots,\allowbreak \alpha_m\) 线性表示，但不能由向量组（I）：\(\displaystyle \alpha_1,\allowbreak \alpha_2,\allowbreak \ldots,\allowbreak \alpha_{m-1}\) 线性表示，记向量组（II）：\(\displaystyle \alpha_1,\allowbreak \alpha_2,\allowbreak \ldots,\allowbreak \alpha_{m-1},\allowbreak \beta\)，则（　）。（A）\(\displaystyle \alpha_m\) 不能由（I）线性表示，也不能由（II）线性表示。（B）\(\displaystyle \alpha_m\) 不能由（I）线性表示，但可由（II）线性表示。（C）\(\displaystyle \alpha_m\) 可由（I）线性表示，也可由（II）线性表示。（D）\(\displaystyle \alpha_m\) 可由（I）线性表示，但不可由（II）线性表示。


（4）设 \(\displaystyle A,\allowbreak B\) 为 \(\displaystyle n\) 阶矩阵，且 \(\displaystyle A\) 与 \(\displaystyle B\) 相似，\(\displaystyle E\) 为 \(\displaystyle n\) 阶单位矩阵，则（　）。


（4）（续）（A）\(\displaystyle \lambda E-A=\lambda E-B\)。（B）\(\displaystyle A\) 与 \(\displaystyle B\) 有相同的特征值和特征向量。（C）\(\displaystyle A\) 与 \(\displaystyle B\) 都相似于一个对角矩阵。（D）对任意常数 \(\displaystyle t\)，\(\displaystyle tE-A\) 与 \(\displaystyle tE-B\) 相似。


（5）设随机变量 \(\displaystyle X_i\sim\begin{pmatrix}-1&0&1\\\dfrac{\displaystyle 1}{\displaystyle 4}&\dfrac{\displaystyle 1}{\displaystyle 2}&\dfrac{\displaystyle 1}{\displaystyle 4}\end{pmatrix}\;(i=1,\allowbreak 2)\)，且满足 \(\displaystyle P\{X_1X_2=0\}=1\)，则 \(\displaystyle P\{X_1=X_2\}\) 等于（　）。（A）\(\displaystyle 0\)。（B）\(\displaystyle \dfrac{\displaystyle 1}{\displaystyle 4}\)。（C）\(\displaystyle \dfrac{\displaystyle 1}{\displaystyle 2}\)。（D）\(\displaystyle 1\)。


\subsection*{三、（本题满分 6 分）}

曲线 \(\displaystyle y=\dfrac{\displaystyle 1}{\displaystyle \sqrt x}\) 的切线与 \(\displaystyle x\) 轴和 \(\displaystyle y\) 轴围成一个图形，记切点的横坐标为 \(\displaystyle a\)。试求切线方程和这个图形的面积。当切点沿曲线趋于无穷远时，该面积的变化趋势如何？


\subsection*{四、（本题满分 7 分）}

计算二重积分 \(\displaystyle \iint_D y\,\mathrm dx\mathrm dy\)，其中 \(\displaystyle D\) 是由直线 \(\displaystyle x=-2,\allowbreak y=0,\allowbreak y=2\) 以及曲线 \(\displaystyle x=-\sqrt{2y-y^2}\) 所围成的平面区域。


\subsection*{五、（本题满分 6 分）}

设生产某种产品必须投入两种要素，\(\displaystyle x_1\) 和 \(\displaystyle x_2\) 分别为两要素的投入量，\(\displaystyle Q\) 为产出量；若生产函数为 \(\displaystyle Q=2x_1^\alpha x_2^\beta\)，其中 \(\displaystyle \alpha,\allowbreak \beta\) 为正常数，且 \(\displaystyle \alpha+\beta=1\)。假设两种要素的价格分别为 \(\displaystyle p_1\) 和 \(\displaystyle p_2\)，试问：当产出量为 12 时，两要素各投入多少可以使得投入总费用最小？


\subsection*{六、（本题满分 6 分）}

设有微分方程 \(\displaystyle y'-2y=\varphi(x)\)，其中 \(\displaystyle \varphi(x)=\begin{cases}2,\allowbreak &x<1,\allowbreak \\0,\allowbreak &x>1.\end{cases}\) 试求在 \(\displaystyle ( -\infty,\allowbreak +\infty )\) 内的连续函数 \(\displaystyle y=y(x)\)，使之在 \(\displaystyle ( -\infty,\allowbreak 1 )\) 和 \(\displaystyle (1,\allowbreak +\infty)\) 内都满足所给方程，且满足条件 \(\displaystyle y(0)=0\)。


\subsection*{七、（本题满分 6 分）}

设函数 \(\displaystyle f(x)\) 连续，且 \(\displaystyle \int_0^x t f(2x-t)\,\mathrm dt=\dfrac{\displaystyle 1}{\displaystyle 2}\arctan x^2\)。已知 \(\displaystyle f(1)=1\)，求 \(\displaystyle \int_1^2 f(x)\,\mathrm dx\) 的值。


\subsection*{八、（本题满分 7 分）}

设函数 \(\displaystyle f(x)\) 在区间 \(\displaystyle [0,\allowbreak 1]\) 上连续，在 \(\displaystyle (0,\allowbreak 1)\) 内可导，且 \(\displaystyle f(0)=f(1)=0\)，\(\displaystyle f(\dfrac{\displaystyle 1}{\displaystyle 2})=1\)。试证：（1）存在 \(\displaystyle \eta\in(\dfrac{\displaystyle 1}{\displaystyle 2},\allowbreak 1)\)，使 \(\displaystyle f(\eta)=\eta\)；（2）对任意实数 \(\displaystyle \lambda\)，必存在 \(\displaystyle \xi\in(0,\allowbreak \eta)\)，使得 \(\displaystyle f'(\xi)-\lambda[f(\xi)-\xi]=1\)。


\subsection*{九、（本题满分 9 分）}

设矩阵 \(\displaystyle A=\begin{pmatrix}a&-1&c\\5&b&3\\1-c&0&-a\end{pmatrix}\)，且 \(\displaystyle |A|=-1\)。又设 \(\displaystyle A\) 的伴随矩阵 \(\displaystyle A^*\) 有特征值 \(\displaystyle \lambda_0\)，属于 \(\displaystyle \lambda_0\) 的特征向量为 \(\displaystyle \alpha=(-1,\allowbreak -1,\allowbreak 1)^{\mathrm T}\)，求 \(\displaystyle a,\allowbreak b,\allowbreak c\) 及 \(\displaystyle \lambda_0\) 的值。


\subsection*{十、（本题满分 7 分）}

设 \(\displaystyle A\) 为 \(\displaystyle m\times n\) 实矩阵，\(\displaystyle E\) 为 \(\displaystyle n\) 阶单位矩阵，已知矩阵 \(\displaystyle B=\lambda E+A^{\mathrm T}A\)，试证：当 \(\displaystyle \lambda>0\) 时，矩阵 \(\displaystyle B\) 为正定矩阵。


\subsection*{十一、（本题满分 9 分）}

假设二维随机变量 \(\displaystyle (X,\allowbreak Y)\) 在矩形 \(\displaystyle G=\{(x,\allowbreak y)\mid 0\le x\le2,\allowbreak 0\le y\le1\}\) 上服从均匀分布，记 \(\displaystyle U=\begin{cases}0,\allowbreak &X\le Y,\allowbreak \\1,\allowbreak &X>Y,\allowbreak \end{cases}\quad V=\begin{cases}0,\allowbreak &X\le2Y,\allowbreak \\1,\allowbreak &X>2Y.\end{cases}\)（1）求 \(\displaystyle U\) 和 \(\displaystyle V\) 的联合分布；（2）求 \(\displaystyle U\) 和 \(\displaystyle V\) 的相关系数 \(\displaystyle r\)。


\subsection*{十二、（本题满分 7 分）}

设 \(\displaystyle X_1,\allowbreak X_2,\allowbreak \ldots,\allowbreak X_9\) 是来自正态总体 \(\displaystyle X\) 的简单随机样本，\(\displaystyle Y_1=\dfrac{\displaystyle 1}{\displaystyle 6}(X_1+\cdots+X_6)\)，\(\displaystyle Y_2=\dfrac{\displaystyle 1}{\displaystyle 3}(X_7+X_8+X_9)\)，\(\displaystyle S^2=\dfrac{\displaystyle 1}{\displaystyle 2}\displaystyle \sum_{i=7}^{9}(X_i-Y_2)^2\)，\(\displaystyle Z=\dfrac{\displaystyle \sqrt2(Y_1-Y_2)}{\displaystyle S}\)，证明统计量 \(\displaystyle Z\) 服从自由度为 2 的 \(\displaystyle t\) 分布。

\end{document}
